\begin{abstract}
O dimensionamento de pontes de madeira envolve múltiplas variáveis de projeto e critérios normativos conflitantes, tradicionalmente conduzido por procedimentos iterativos dependentes da experiência do projetista, sem investigação sistemática de como vão, classe de resistência e variabilidade dimensional interagem no comportamento estrutural.
Este trabalho propõe um framework paramétrico de otimização robusta multiobjetivo para pontes de madeira roliça, com duas funções objetivo antagônicas, o volume de madeira e a utilização do limite de deslocamento, resolvidas pelo NSGA-II sob as restrições de estado limite último e de serviço da ABNT NBR 7190:2022.
A robustez é incorporada mantendo os objetivos no ponto nominal e agregando as restrições pelo pior caso sobre uma grade determinística de cinco pontos de desvio $\pm\rho$ ($\rho = 5\%$) no diâmetro da longarina.
O framework foi aplicado a 20 configurações, combinando quatro vãos (3{,}0 a 6{,}0~m) e as cinco classes de resistência D20 a D60 sob o trem tipo TB-240.
Todas convergiram para soluções viáveis, com consumo de madeira entre 1{,}75 e 9{,}98~m\textsuperscript{3}.
O consumo cresce com o vão segundo uma lei de potência de expoente entre 1{,}64 e 1{,}73, e o primeiro degrau de classe, D20 para D30, responde por 41 a 48\% de toda a economia até D60.
A flexão governa a totalidade das células, quinze pelo tabuleiro e cinco pela longarina, sem nenhuma governada pela flecha, que ainda assim se aproxima do limite nas células mais severas.
A robustez tem custo real e monotônico, de $+6{,}0\%$ ($\rho=5\%$) a $+26{,}9\%$ ($\rho=20\%$) em volume.
Os resultados foram condensados em ábacos de anteprojeto e em uma equação de pré-dimensionamento ($R^2=0{,}999$), válida no domínio investigado, quantificando como vão, classe e variabilidade dimensional controlam a demanda de material e os estados limites governantes.
\end{abstract}


\textbf{Palavras-chave}: pontes de madeira roliça, otimização multiobjetivo, projeto robusto, projeto paramétrico, otimização estrutural, estruturas de madeira.
